\paragraph{An archimedean estimate on each complex theta fibre.}
Fix a complex embedding \(\sigma\) of the coefficient field \(F\),
of degree \(h\) over \(\mathbf Q\). Use the cubical metrics
rigidified to have norm one at zero on the complex graph lines and
Haar probability on each abelian variety. Put \(g=fr=\dim A\),
\(m=gn\), \(d=\deg P\), and \(\Delta=|\operatorname{disc}E|\).
Let \(\lambda>0\) be the shortest nonzero period of \(A\) in the
positive Hermitian metric \(H\) of its principal polarization,
normalized by the Riemann form \(\Im H\). Take an orthogonal Hodge
eigenframe whose individual direction vectors have length at most
\(U\), with \(U>0\). Use the row copies of that frame in \(A^n\).
Let \(\|P^\sigma\|_1\) denote the sum of absolute values of its
coefficients in these directions. The integer \(N\) and the self-dual
Hermitian column lattice \(\Lambda\) can vary at this fixed frame.

\begin{proposition}[A complex fibre bound independent of the auxiliary level]
\label{intjet:arch}
With the rescaling \(s\) chosen above by Minkowski's theorem, put
\[
 C=2fn\Delta,\qquad
 R=\frac{\min(1,\lambda/C)}{4\sqrt m}.
\]
Then the cubically normalized value of the rescaled theta series satisfies
\[
 \frac{|\sigma(sy)(x_0)|}{\|s_0\|_2^n}
 \le |\Delta_F|^{1/(2h)}e^{\pi/2}
       \|P^\sigma\|_1\,d!\,(\sqrt2\,U)^d R^{-(d+m)}.
\]
This estimate is uniform in \(N,\Lambda\), the choice of division
points and the level character.
\end{proposition}
\begin{proof}
Use the isogeny \(a:B=A^n\to B_1=\Lambda\otimes_O A\) from
Theorem~\ref{intjet:trace}; in rational column coordinates its
differential is \(N\) times the identity. Both graph lines on
\(B\) and \(B_1\) induce principal polarizations. For the latter
this follows by tensoring the principal polarization of \(A\) with
the isomorphism from \(\Lambda\) to its conjugate \(O\)-dual
given by the unimodular Hermitian pairing. The equality
\(a^*\mathcal L_1=[N]^*\mathcal L\) is isometric: the two
curvature forms agree and the ratio of metrics on a compact complex
torus is then constant, fixed at zero. A finite-order normalized
translation lift is also an isometry, since the ratio is constant
and a positive real number of finite multiplicative order is one.
All the \(L^2\) pullbacks under isogenies preserve norms because
the pushforward of Haar probability is Haar probability.

Write \(v=[N]^*s_0^{\boxtimes n}\), and \(U_z\) for the normalized
translation. The vectors \(U_zv\), \(z\in Z_\Lambda/K_0\), are
pairwise orthogonal. Indeed \(v\) is fixed by the deck translations
\(b\in B[N]\), and \(U_zv\) is then an eigenvector for each
\(U_b\), with character given by the theta commutator with \(z\).
The annihilator of \(B[N]\) inside \(B[N^2]\) for this commutator
is \(B[N]\) itself. In a symplectic basis for the principal
polarization, the pairing is the standard perfect alternating
pairing modulo \(N^2\), and \(B[N]\) consists of all multiples of
\(N\). Distinct classes modulo \(K_0=Z_\Lambda\cap B[N]\) therefore
give distinct characters of commuting unitaries. The trace identity
of Theorem~\ref{intjet:trace} gives
\[
 a^*T=\sum_{z\in Z_\Lambda/K_0}U_zv,\qquad
 \|T\|_2= I^{r/2}\|s_0\|_2^n,\qquad I=[\Lambda:\Lambda_0].
\]
In particular \(\|\sigma(s)T\|_2\le
|\Delta_F|^{1/(2h)}\|s_0\|_2^n\). This equality of squared norms
uses the original order of the translations. The \(t_u\) act
after this calculation.

We bound the shortest period on \(B_1\). On \(\Lambda\) take the
real trace norm \(\|v\|_{\rm tr}^2=
\operatorname{Tr}_{E/\mathbf Q}\langle v,v\rangle\).
For \(v\ne0\) its squared norm is at least one, since the trace
is a positive rational integer. The trace-dual lattice equals
\(\mathfrak D_E^{-1}\Lambda\), by \(O\)-duality and the definition
of the codifferent. Multiplication by \(\Delta\) sends it into
\(\Lambda\), so its minimum is at least \(1/\Delta\).
Euclidean transference in dimension \(2fn\) gives \(2fn\) independent
vectors of \(\Lambda\) of trace norm at most \(C\)
\cite[Theorem 2.1]{banaszczyk}. Contractions with those vectors
separate every nonzero \(x\in\Lambda\otimes_O\Pi\), where
\(\Pi\) is the period lattice of \(A\): their rational span is the
whole column space over \(\mathbf Q\), and the pairing is perfect
over \(E\). Each contraction belongs to \(\Pi\). Its operator norm
on the complex tangent space with product Hermitian norm is at
most \(C\), by Cauchy--Schwarz on each eigenblock; on a given block
the sum of the squared column coefficients is at most the squared
trace norm. The positive form on \(B_1\) in rational column
coordinates is this same product form by the isometry of pullbacks
and \(da=N\,\mathrm{id}\). Thus every nonzero period on \(B_1\)
has norm at least \(\lambda/C\).

Use orthonormal eigen-coordinates on the universal cover of \(B_1\)
centered at any prescribed point. The polydisc of coordinate radius
\(R\) injects, by the period bound. In it a holomorphic line frame
of squared norm \(\exp(-\pi\|z\|^2)\) exists and is unique up to
unit phase. Indeed its prescribed logarithmic norm has the required
curvature, and a real pluriharmonic discrepancy on the simply
connected chart is the real part of a holomorphic function.
The Euclidean covolume of the periods in these coordinates is one:
the real volume form is \((\Im H_1)^m/m!\) up to orientation,
and its period integral for a principal polarization is one.
The same holds for any translate of the line and any choice of
center. If \(G=\sum c_\alpha z^\alpha\) is the coefficient of a
section of \(L^2\)-norm \(D\) in such a frame, then
\[
 |c_\alpha|^2\,
 \frac{\pi^mR^{2|\alpha|+2m}}{\prod_i(\alpha_i+1)}
 \le \int_{\{|z_i|<R\}}|G|^2\,dz\le e^\pi D^2 .
\]
The first inequality is angular orthogonality followed by radial
integration. The second uses \(mR^2\le1\) and injectivity. Since
\(\prod(\alpha_i+1)\le2^{|\alpha|}\), derivatives of degree \(d\)
in the row frame are bounded by
\(e^{\pi/2}D\,d!(\sqrt2 U)^d R^{-d-m}\).
Here each row direction is a multiple of an orthonormal coordinate
direction. Summing with \(P\) introduces \(\|P^\sigma\|_1\).

To compare this local derivative with the harmonic jet, on the
vector-space cover choose the holomorphic frame of the polarized
biextension with squared norm
\(\exp[-\pi(H(z,w)+H(w,z))]\), normalized at the origin.
Its curvature is that of the polarized biextension, by taking the
pullback under addition and dividing by the two projections.
Addition in either variable and Rosati adjunction preserve this
frame: their metric factors agree by sesquilinearity, and the ratio
is holomorphic of unit modulus and is one on the axes. On the graph
\(w=w_0z\) its squared norm is \(\exp(-\pi H(z,z))\) since
\(w_0+w_0^\dagger=1\). Make this frame comparison first on the
single factor \(A\times A\). Any permitted additive frame there
differs by the exponential of a Rosati-compatible complex bilinear
form, hence only couples conjugate types. On \(B=A^n\) the jet
uses the exterior product of the same single-factor frame in all
columns, followed by its \(N^2\)-th tensor power. Its graph correction
is therefore a sum of scalar multiples of
\(\sum_b z_{\tau i b}z_{\bar\tau j b}\). Each such sum disappears
under the harmonic equations, as in Lemma~\ref{intjet:der}.
For the final translation
by \(t_u\), apply the centered frame at \(a(t_u)\) to the translated
germ. Its pullback has squared frame norm
\(\exp(-\pi N^2\|z\|^2)\); the isometry \(U_{t_u}\) identifies it
with the frame at zero up to unit phase, by uniqueness. The derivative
then gains precisely \(N^d\), which cancels the normalizing
denominator in the series. Finally average over \(u\). Its Fourier
weights have modulus one and the divisor is their number \(C_M\).
This proves the inequality.
\end{proof}

Write \(S=d+m/2\) and \(Z_0=frL_0\). Suppose the frame satisfies
\[
 \max\{\log^+(\lambda^{-1}),\log^+ U,\log^+\|s_0\|_2\}\ll Z_0^2,
 \qquad \log^+\|P^\sigma\|_1\ll Z_0^2S .
\]
Then the log of the absolute bound is \(O(Z_0^2 S)\) under the
stage bounds above.
Here \(\log^+x=\max(0,\log x)\) for \(x>0\); a zero polynomial is
immediate. In fact \(m\le2S\), \(d\le S\),
\(\log^+(d!)\le d\log(1+d)\), \(\log(2S)\ll L_0\),
\(\log\Delta\le7fL_0\) and \(h^{-1}\log|\Delta_F|\ll Z_0^2\).
These bound every factor in the displayed inequality, including
\(\log(1/R)\ll Z_0^2\). To pass from this complex fibre estimate
to the simultaneous section norms requires the normalized algebraic
and analytic readings on every component, short polynomials spanning
the actual space, and the indicated frame control at its half
maxima. The normalized norm of the chosen generator matters when
that generator is identified with a frame of the determinant line.
\src{src/Development/AbelianVarietyTheory/Proofs/CuspFormEvaluations/NormalisedThetaValueBound.lean}{146}{Normalization on the tensor torus};
\src{src/Development/AbelianVarietyTheory/Proofs/DivisorProofs/NormRawThetaSeriesIndexBoundLeSqrt.lean}{18}{Square-root index in the archimedean estimate}.
